\subsection{The field of fractions of the ring of formal power series in one indeterminate over a field}

If K is a commutative field, we shall denote by \(\mathrm{K}((\mathbf{X}))\) the field of fractions of the integral domain \(\mathrm{K}[[\mathbf{X}]]\) .

PROPOSITION 12. --- Every non-zero element u of K((X)) may be written in a unique way as \(u = X^{k}v\) , where \(k \in Z\) and v is a formal power series in X of order 0. Let u = w/t, where w, t are non-zero elements of K{[}{[}X{]}{]}. We have \(w = X^{r}w_{1}, t = X^{s}t_{1}\) , where r, \(s \in N\) and \(w_{1}, t_{1}\) are formal power series of order 0, hence invertible in K{[}{[}X{]}{]} (IV, p.~30, Prop. 6). Then \(u = X^{r-s}w_{1}t_{1}^{-1}\) and \(w_{1}t_{1}^{-1}\) is a formal power series of order 0.

Let us prove the uniqueness. Suppose that \(u = X^{k_{1}}v_{1} = X^{k_{2}}v_{2}\) where \(k_{1}, k_{2} \in Z\) and \(v_{1}, v_{2}\) are formal power series of order 0. Since \(X^{k_{1}-k_{2}} = v_{2}v_{1}^{-1}\) is a formal power series of order 0, we have \(k_{1} = k_{2}\) whence \(v_{1} = v_{2}\) and this proves the uniqueness assertion.

We shall say that the elements of \(\mathbf{K}((\mathbf{X}))\) are generalized formal power series in X with coefficients in K, or simply formal power series when no confusion can arise (the elements of \(K[[X]]\) are then called formal power series with positive exponents); if \(u \neq 0\) , the integer k defined in Prop. 12 is also called the order of u and is written \(\omega(u)\) , even if it is \textless{} 0; we also put \(\omega(0) = \infty\) . It may be verified at once that

\[
\begin{array}{l} \omega (u + v) \geqslant \inf (\omega (u), \omega (v)) \\ \omega (u + v) = \inf (\omega (u), \omega (v)) \quad \text { if } \quad \omega (u) \neq \omega (v) \\ \omega (u v) = \omega (u) + \omega (v) \end{array}
\]

still hold for generalized formal power series. In particular, if \(u \neq 0\) , then \(\omega(u^{-1}) = -\omega(u)\) . * In other words (Comm. Alg., VI, § 3, No.~6, p.~392, Def. 3), \(\omega\) is a normalized discrete valuation of the field \(K((X))\) . *

For each integer \(n \in \mathbf{Z}\) let \(\mathfrak{p}_n\) be the set of all \(u \in K((X))\) such that \(\omega(u) \geqslant n\) . Then \((\mathfrak{p}_n)_{n \in \mathbf{Z}}\) is a decreasing sequence of subgroups of the additive group \(K((X))\) , with intersection 0; there exists thus a topology on \(K((X))\) , invariant under translation, for which \((\mathfrak{p}_n)_{n \in \mathbf{Z}}\) is a fundamental system of neighbourhoods of 0 (Gen.~Top., III, p.~223). We can easily verify that \(K((X))\) is a topological field (Gen.~Top., III, p.~281) and that \(K[[X]]\) is an open and closed subspace of \(K((X))\) .

Let \((\alpha_{n})_{n\in \mathbf{Z}}\) be a family of elements of \(\mathbf{K}\) , and suppose that there exists an integer \(\mathbf{N}\) such that \(\alpha_{n} = 0\) for all \(n < \mathbf{N}\) . Then the family \((\alpha_{n}\mathbf{X}^{n})_{n\in \mathbf{Z}}\) is summable in \(\mathbf{K}((\mathbf{X}))\) (Gen.~Top., III, p.~263, Cor.); put \(u = \sum_{n\in \mathbf{Z}}\alpha_n\mathbf{X}^n\) , then \(u = 0\) if and

only if \(\alpha_{n} = 0\) for all \(n\) ; otherwise the order of \(u\) is the least integer \(k\) such that \(\alpha_{k} \neq 0\) . Finally every element of \(\mathbf{K}((\mathbf{X}))\) may be written in a unique fashion in the form \(\sum_{n \in \mathbf{Z}} \alpha_{n} \mathbf{X}^{n}\) , where the sequence \((\alpha_{n})\) satisfies \(\alpha_{-n} = 0\) for all sufficiently large \(n\).

Since the ring K{[}X{]} is a subring of K{[}{[}X{]}{]}, every rational fraction \(u/v \in \mathrm{K}(\mathrm{X})\) (u, v being polynomials in X) may be identified with the (generalized) formal power series \(uv^{-1}\) of K((X)), which we shall call its expansion at the origin; the field K(X) is thus identified with a subfield of K((X)).

